\documentclass[12pt]{report}
\usepackage[utf8x]{inputenc}

%Warning: tipa declares many non-standard macros used by utf8x to
%interpret utf8 characters but extra packages might have to be added
%such as "textgreek" for Greek letters not already in tipa
%or "stmaryrd" for mathematical symbols.
%Utf8 codes missing a LaTeX interpretation can be defined by using
%\DeclareUnicodeCharacter{code}{interpretation}.
%Use coqdoc's option -p to add new packages or declarations.
\usepackage{tipa}

\usepackage[T1]{fontenc}
\usepackage{fullpage}
\usepackage{coqdoc}
\usepackage{amsmath,amssymb}
\usepackage{url}
\begin{document}
%%%%%%%%%%%%%%%%%%%%%%%%%%%%%%%%%%%%%%%%%%%%%%%%%%%%%%%%%%%%%%%%%
%% This file has been automatically generated with the command
%% /nix/store/5nhkcrnn3hzlx1rsgy6zxlg2lrla4f5y-coq-with-packages-9.1.1/bin-orig/coqdoc -latex -toc -interpolate -utf8 -Q . SECF Maps.v Imp.v Equiv.v Hoare.v Hoare2.v Preface.v Noninterference.v StaticIFC.v SpecCT.v Postscript.v Bib.v MapsTest.v ImpTest.v EquivTest.v HoareTest.v Hoare2Test.v PrefaceTest.v NoninterferenceTest.v StaticIFCTest.v SpecCTTest.v PostscriptTest.v BibTest.v 
%%%%%%%%%%%%%%%%%%%%%%%%%%%%%%%%%%%%%%%%%%%%%%%%%%%%%%%%%%%%%%%%%
\coqlibrary{SECF.Maps}{Library }{SECF.Maps}

\begin{coqdoccode}
\end{coqdoccode}
\section{Maps: Total and Partial Maps}



 \textit{Maps} (or \textit{dictionaries}) are ubiquitous data structures both in
    ordinary programming and in the theory of programming languages;
    we're going to need them in many places in the coming chapters.


    They also make a nice case study using ideas we've seen in
    previous chapters, including building data structures out of
    higher-order functions (from \coqdocvar{Basics} and \coqdocvar{Poly}) and the use of
    reflection to streamline proofs (from \coqdocvar{IndProp}).


    We'll define two flavors of maps: \textit{total} maps, which include a
    ``default'' element to be returned when a key being looked up
    doesn't exist, and \textit{partial} maps, which instead return an
    \coqexternalref{option}{http://rocq-prover.org/doc/V9.1.1/stdlib//Corelib.Init.Datatypes}{\coqdocinductive{option}} to indicate success or failure.  Partial maps are defined
    in terms of total maps, using \coqexternalref{None}{http://rocq-prover.org/doc/V9.1.1/stdlib//Corelib.Init.Datatypes}{\coqdocconstructor{None}} as the default element. 
\begin{coqdoccode}
\coqdocemptyline
\end{coqdoccode}
\section{The  Standard Library}



 One small digression before we begin...


    Unlike the chapters we have seen so far, this one does not
    \coqdockw{Require} \coqdockw{Import} the chapter before it (or, transitively, all the
    earlier chapters).  Instead, in this chapter and from now on,
    we're going to import the definitions and theorems we need
    directly from Rocq's standard library.  You should not notice much
    difference, though, because we've been careful to name our own
    definitions and theorems the same as their counterparts in the
    standard library, wherever they overlap. 
\begin{coqdoccode}
\coqdocemptyline
\coqdocnoindent
\coqdockw{From} \coqdocvar{Stdlib} \coqdockw{Require} \coqdockw{Import} \coqdoclibrary{Arith}.\coqdoceol
\coqdocnoindent
\coqdockw{From} \coqdocvar{Stdlib} \coqdockw{Require} \coqdockw{Import} \coqdoclibrary{Bool}.\coqdoceol
\coqdocnoindent
\coqdockw{From} \coqdocvar{Stdlib} \coqdockw{Require} \coqdockw{Export} \coqdoclibrary{Strings.String}.\coqdoceol
\coqdocnoindent
\coqdockw{From} \coqdocvar{Stdlib} \coqdockw{Require} \coqdockw{Import} \coqdoclibrary{FunctionalExtensionality}.\coqdoceol
\coqdocnoindent
\coqdockw{From} \coqdocvar{Stdlib} \coqdockw{Require} \coqdockw{Import} \coqdoclibrary{List}.\coqdoceol
\coqdocnoindent
\coqdockw{Import} \coqdocmodule{ListNotations}.\coqdoceol
\coqdocemptyline
\end{coqdoccode}
Documentation for the standard library can be found at
    https://rocq-prover.org/doc/V9.0.0/stdlib/index.html.


    The \coqdockw{Search} command is a good way to look for theorems involving
    objects of specific types. See \coqdocvar{Lists} for a reminder of how
    to use it. 

 If you want to find out how or where a notation is defined, the
    \coqdockw{Locate} command is useful.  For example, where is the natural
    addition operation defined in the standard library? 
\begin{coqdoccode}
\coqdocemptyline
\coqdocnoindent
\coqdockw{Locate} "+".\coqdoceol
\coqdocemptyline
\end{coqdoccode}
(There are several uses of the + notation, but only one for
    naturals.) 
\begin{coqdoccode}
\coqdocemptyline
\coqdocnoindent
\coqdockw{Print} \coqdocvar{Init.Nat.add}.\coqdoceol
\coqdocemptyline
\end{coqdoccode}
We'll see some more uses of \coqdockw{Locate} in the \coqdoclibrary{Imp} chapter. 
\begin{coqdoccode}
\coqdocemptyline
\end{coqdoccode}
\section{Identifiers}



 To define maps, we first need a type for the keys that we will use
    to index into our maps.  In \coqdocvar{Lists.v} we introduced a fresh type
    \coqdocvar{id} for a similar purpose; here and for the rest of \textit{Software
    Foundations} we will use the \coqdocinductive{string} type from Rocq's standard
    library. 

 To compare strings, we use the function \coqdoclemma{eqb\_refl} from the \coqdoclibrary{String}
    module in the standard library. 
\begin{coqdoccode}
\coqdocemptyline
\coqdocnoindent
\coqdockw{Check} \coqdoclemma{String.eqb\_refl} :\coqdoceol
\coqdocindent{1.00em}
\coqdockw{\ensuremath{\forall}} \coqdef{SECF.Maps.x:1}{x}{\coqdocbinder{x}} : \coqdocinductive{string}, (\coqref{SECF.Maps.x:1}{\coqdocvariable{x}} \coqdocnotation{=?} \coqref{SECF.Maps.x:1}{\coqdocvariable{x}})\%\coqdocvar{string} \coqexternalref{::type scope:x '=' x}{http://rocq-prover.org/doc/V9.1.1/stdlib//Corelib.Init.Logic}{\coqdocnotation{=}} \coqexternalref{true}{http://rocq-prover.org/doc/V9.1.1/stdlib//Corelib.Init.Datatypes}{\coqdocconstructor{true}}.\coqdoceol
\coqdocemptyline
\end{coqdoccode}
We will often use a few basic properties of string equality... 
\begin{coqdoccode}
\coqdocnoindent
\coqdockw{Check} \coqdoclemma{String.eqb\_eq} :\coqdoceol
\coqdocindent{1.00em}
\coqdockw{\ensuremath{\forall}} \coqdef{SECF.Maps.n:2}{n}{\coqdocbinder{n}} \coqdef{SECF.Maps.m:3}{m}{\coqdocbinder{m}} : \coqdocinductive{string}, (\coqref{SECF.Maps.n:2}{\coqdocvariable{n}} \coqdocnotation{=?} \coqref{SECF.Maps.m:3}{\coqdocvariable{m}})\%\coqdocvar{string} \coqexternalref{::type scope:x '=' x}{http://rocq-prover.org/doc/V9.1.1/stdlib//Corelib.Init.Logic}{\coqdocnotation{=}} \coqexternalref{true}{http://rocq-prover.org/doc/V9.1.1/stdlib//Corelib.Init.Datatypes}{\coqdocconstructor{true}} \coqexternalref{::type scope:x '<->' x}{http://rocq-prover.org/doc/V9.1.1/stdlib//Corelib.Init.Logic}{\coqdocnotation{\ensuremath{\leftrightarrow}}} \coqref{SECF.Maps.n:2}{\coqdocvariable{n}} \coqexternalref{::type scope:x '=' x}{http://rocq-prover.org/doc/V9.1.1/stdlib//Corelib.Init.Logic}{\coqdocnotation{=}} \coqref{SECF.Maps.m:3}{\coqdocvariable{m}}.\coqdoceol
\coqdocnoindent
\coqdockw{Check} \coqdoclemma{String.eqb\_neq} :\coqdoceol
\coqdocindent{1.00em}
\coqdockw{\ensuremath{\forall}} \coqdef{SECF.Maps.n:4}{n}{\coqdocbinder{n}} \coqdef{SECF.Maps.m:5}{m}{\coqdocbinder{m}} : \coqdocinductive{string}, (\coqref{SECF.Maps.n:4}{\coqdocvariable{n}} \coqdocnotation{=?} \coqref{SECF.Maps.m:5}{\coqdocvariable{m}})\%\coqdocvar{string} \coqexternalref{::type scope:x '=' x}{http://rocq-prover.org/doc/V9.1.1/stdlib//Corelib.Init.Logic}{\coqdocnotation{=}} \coqexternalref{false}{http://rocq-prover.org/doc/V9.1.1/stdlib//Corelib.Init.Datatypes}{\coqdocconstructor{false}} \coqexternalref{::type scope:x '<->' x}{http://rocq-prover.org/doc/V9.1.1/stdlib//Corelib.Init.Logic}{\coqdocnotation{\ensuremath{\leftrightarrow}}} \coqref{SECF.Maps.n:4}{\coqdocvariable{n}} \coqexternalref{::type scope:x '<>' x}{http://rocq-prover.org/doc/V9.1.1/stdlib//Corelib.Init.Logic}{\coqdocnotation{\ensuremath{\not=}}} \coqref{SECF.Maps.m:5}{\coqdocvariable{m}}.\coqdoceol
\coqdocnoindent
\coqdockw{Check} \coqdoclemma{String.eqb\_spec} :\coqdoceol
\coqdocindent{1.00em}
\coqdockw{\ensuremath{\forall}} \coqdef{SECF.Maps.x:6}{x}{\coqdocbinder{x}} \coqdef{SECF.Maps.y:7}{y}{\coqdocbinder{y}} : \coqdocinductive{string}, \coqexternalref{reflect}{http://rocq-prover.org/doc/V9.1.1/stdlib//Corelib.Init.Datatypes}{\coqdocinductive{reflect}} (\coqref{SECF.Maps.x:6}{\coqdocvariable{x}} \coqexternalref{::type scope:x '=' x}{http://rocq-prover.org/doc/V9.1.1/stdlib//Corelib.Init.Logic}{\coqdocnotation{=}} \coqref{SECF.Maps.y:7}{\coqdocvariable{y}}) (\coqdocdefinition{String.eqb} \coqref{SECF.Maps.x:6}{\coqdocvariable{x}} \coqref{SECF.Maps.y:7}{\coqdocvariable{y}}).\coqdoceol
\coqdocemptyline
\end{coqdoccode}
\section{Total Maps}



 Our main job in this chapter will be to build a definition of
    partial maps that is similar in behavior to the one we saw in the
    \coqdocvar{Lists} chapter, plus accompanying lemmas about its behavior.


    This time around, though, we're going to use \textit{functions}, rather
    than lists of key-value pairs, to build maps.  The advantage of
    this representation is that it offers a more ``extensional'' view of
    maps: two maps that respond to queries in the same way will be
    represented as exactly the same function, rather than just as
    ``equivalent'' list structures.  This simplifies proofs that use
    maps. 

 We build up to partial maps in two steps.  First, we define a type
    of \textit{total maps} that return a default value when we look up a key
    that is not present in the map. 
\begin{coqdoccode}
\coqdocemptyline
\coqdocnoindent
\coqdockw{Definition} \coqdef{SECF.Maps.total map}{total\_map}{\coqdocdefinition{total\_map}} (\coqdef{SECF.Maps.A:8}{A}{\coqdocbinder{A}} : \coqdockw{Type}) := \coqdocinductive{string} \coqexternalref{::type scope:x '->' x}{http://rocq-prover.org/doc/V9.1.1/stdlib//Corelib.Init.Logic}{\coqdocnotation{\ensuremath{\rightarrow}}} \coqref{SECF.Maps.A:8}{\coqdocvariable{A}}.\coqdoceol
\coqdocemptyline
\end{coqdoccode}
Intuitively, a total map over an element type \coqdocvar{A} is just a
    function that can be used to look up \coqdocinductive{string}s, yielding \coqdocvar{A}s. 

 The function \coqref{SECF.Maps.t empty}{\coqdocdefinition{t\_empty}} yields an empty total map, given a default
    element; this map always returns the default element when applied
    to any string. 
\begin{coqdoccode}
\coqdocemptyline
\coqdocnoindent
\coqdockw{Definition} \coqdef{SECF.Maps.t empty}{t\_empty}{\coqdocdefinition{t\_empty}} \{\coqdef{SECF.Maps.A:9}{A}{\coqdocbinder{A}} : \coqdockw{Type}\} (\coqdef{SECF.Maps.v:10}{v}{\coqdocbinder{v}} : \coqref{SECF.Maps.A:9}{\coqdocvariable{A}}) : \coqref{SECF.Maps.total map}{\coqdocdefinition{total\_map}} \coqref{SECF.Maps.A:9}{\coqdocvariable{A}} :=\coqdoceol
\coqdocindent{1.00em}
(\coqdockw{fun} \coqdocvar{\_} \ensuremath{\Rightarrow} \coqref{SECF.Maps.v:10}{\coqdocvariable{v}}).\coqdoceol
\coqdocemptyline
\end{coqdoccode}
More interesting is the map-updating function, which (as always)
    takes a map \coqdocvar{m}, a key \coqdocvar{x}, and a value \coqdocvar{v} and returns a new map
    that takes \coqdocvar{x} to \coqdocvar{v} and takes every other key to whatever \coqdocvar{m}
    does.  The novelty here is that we achieve this effect by wrapping
    a new function around the old one. 
\begin{coqdoccode}
\coqdocemptyline
\coqdocnoindent
\coqdockw{Definition} \coqdef{SECF.Maps.t update}{t\_update}{\coqdocdefinition{t\_update}} \{\coqdef{SECF.Maps.A:11}{A}{\coqdocbinder{A}} : \coqdockw{Type}\} (\coqdef{SECF.Maps.m:12}{m}{\coqdocbinder{m}} : \coqref{SECF.Maps.total map}{\coqdocdefinition{total\_map}} \coqref{SECF.Maps.A:11}{\coqdocvariable{A}})\coqdoceol
\coqdocindent{10.00em}
(\coqdef{SECF.Maps.x:13}{x}{\coqdocbinder{x}} : \coqdocinductive{string}) (\coqdef{SECF.Maps.v:14}{v}{\coqdocbinder{v}} : \coqref{SECF.Maps.A:11}{\coqdocvariable{A}}) :=\coqdoceol
\coqdocindent{1.00em}
\coqdockw{fun} \coqdef{SECF.Maps.x':15}{x'}{\coqdocbinder{x'}} \ensuremath{\Rightarrow} \coqdockw{if} \coqdocdefinition{String.eqb} \coqref{SECF.Maps.x:13}{\coqdocvariable{x}} \coqref{SECF.Maps.x':15}{\coqdocvariable{x'}} \coqdockw{then} \coqref{SECF.Maps.v:14}{\coqdocvariable{v}} \coqdockw{else} \coqref{SECF.Maps.m:12}{\coqdocvariable{m}} \coqref{SECF.Maps.x':15}{\coqdocvariable{x'}}.\coqdoceol
\coqdocemptyline
\end{coqdoccode}
This definition is a nice example of higher-order programming:
    \coqref{SECF.Maps.t update}{\coqdocdefinition{t\_update}} takes a \textit{function} \coqdocvar{m} and yields a new function
    \coqdockw{fun} \coqdocvar{x'} \ensuremath{\Rightarrow} ... that behaves like the desired map. 

 For example, we can build a map taking \coqdocinductive{string}s to \coqexternalref{bool}{http://rocq-prover.org/doc/V9.1.1/stdlib//Corelib.Init.Datatypes}{\coqdocinductive{bool}}s, where
    "\coqref{SECF.Imp.AExp.foo}{\coqdoclemma{foo}}" and "\coqdocvar{bar}" are mapped to \coqexternalref{true}{http://rocq-prover.org/doc/V9.1.1/stdlib//Corelib.Init.Datatypes}{\coqdocconstructor{true}} and every other key is
    mapped to \coqexternalref{false}{http://rocq-prover.org/doc/V9.1.1/stdlib//Corelib.Init.Datatypes}{\coqdocconstructor{false}}, like this: 
\begin{coqdoccode}
\coqdocemptyline
\coqdocnoindent
\coqdockw{Definition} \coqdef{SECF.Maps.examplemap}{examplemap}{\coqdocdefinition{examplemap}} :=\coqdoceol
\coqdocindent{1.00em}
\coqref{SECF.Maps.t update}{\coqdocdefinition{t\_update}} (\coqref{SECF.Maps.t update}{\coqdocdefinition{t\_update}} (\coqref{SECF.Maps.t empty}{\coqdocdefinition{t\_empty}} \coqexternalref{false}{http://rocq-prover.org/doc/V9.1.1/stdlib//Corelib.Init.Datatypes}{\coqdocconstructor{false}}) "foo" \coqexternalref{true}{http://rocq-prover.org/doc/V9.1.1/stdlib//Corelib.Init.Datatypes}{\coqdocconstructor{true}})\coqdoceol
\coqdocindent{5.50em}
"bar" \coqexternalref{true}{http://rocq-prover.org/doc/V9.1.1/stdlib//Corelib.Init.Datatypes}{\coqdocconstructor{true}}.\coqdoceol
\coqdocemptyline
\end{coqdoccode}
Next, let's introduce some notations to facilitate working with
    maps. 

 First, we use the following notation to represent an empty total
    map with a default value. 
\begin{coqdoccode}
\coqdocnoindent
\coqdockw{Notation} \coqdef{SECF.Maps.:::'  ' '!->' x}{"}{"}'\_\_' '!->' v" := (\coqref{SECF.Maps.t empty}{\coqdocdefinition{t\_empty}} \coqdocvar{v})\coqdoceol
\coqdocindent{1.00em}
(\coqdoctac{at} \coqdockw{level} 100, \coqdoctac{right} \coqdockw{associativity}).\coqdoceol
\coqdocemptyline
\coqdocnoindent
\coqdockw{Example} \coqdef{SECF.Maps.example empty}{example\_empty}{\coqdocdefinition{example\_empty}} := ( \coqexternalref{false}{http://rocq-prover.org/doc/V9.1.1/stdlib//Corelib.Init.Datatypes}{\coqdocconstructor{false}}).\coqdoceol
\coqdocemptyline
\end{coqdoccode}
We next introduce a symbolic notation for extending an existing
    map with a new binding. 
\begin{coqdoccode}
\coqdocnoindent
\coqdockw{Notation} \coqdef{SECF.Maps.:::x '!->' x ';' x}{"}{"}x '!->' v ';' m" := (\coqref{SECF.Maps.t update}{\coqdocdefinition{t\_update}} \coqdocvar{m} \coqdocvar{x} \coqdocvar{v})\coqdoceol
\coqdocindent{16.00em}
(\coqdoctac{at} \coqdockw{level} 100, \coqdocvar{v} \coqdockw{constr} \coqdoctac{at} \coqdockw{level} 100, \coqdoctac{right} \coqdockw{associativity}).\coqdoceol
\coqdocemptyline
\end{coqdoccode}
The \coqref{SECF.Maps.examplemap}{\coqdocdefinition{examplemap}} above can now be defined as follows: 
\begin{coqdoccode}
\coqdocemptyline
\coqdocnoindent
\coqdockw{Definition} \coqdef{SECF.Maps.examplemap'}{examplemap'}{\coqdocdefinition{examplemap'}} :=\coqdoceol
\coqdocindent{1.00em}
( "bar" \coqref{SECF.Maps.:::x '!->' x ';' x}{\coqdocnotation{!->}} \coqexternalref{true}{http://rocq-prover.org/doc/V9.1.1/stdlib//Corelib.Init.Datatypes}{\coqdocconstructor{true}}\coqref{SECF.Maps.:::x '!->' x ';' x}{\coqdocnotation{;}}\coqdoceol
\coqdocindent{2.00em}
"foo" \coqref{SECF.Maps.:::x '!->' x ';' x}{\coqdocnotation{!->}} \coqexternalref{true}{http://rocq-prover.org/doc/V9.1.1/stdlib//Corelib.Init.Datatypes}{\coqdocconstructor{true}}\coqref{SECF.Maps.:::x '!->' x ';' x}{\coqdocnotation{;}}\coqdoceol
\coqdocindent{2.00em}
\coqref{SECF.Maps.:::'  ' '!->' x}{\coqdocnotation{\_\_}}    \coqref{SECF.Maps.:::'  ' '!->' x}{\coqdocnotation{!->}} \coqexternalref{false}{http://rocq-prover.org/doc/V9.1.1/stdlib//Corelib.Init.Datatypes}{\coqdocconstructor{false}}\coqdoceol
\coqdocindent{1.00em}
).\coqdoceol
\coqdocemptyline
\end{coqdoccode}
This completes the definition of total maps.  Note that we
    don't need to define a \coqdocvar{find} operation on this representation of
    maps because it is just function application! 
\begin{coqdoccode}
\coqdocemptyline
\coqdocnoindent
\coqdockw{Example} \coqdef{SECF.Maps.update example1}{update\_example1}{\coqdocdefinition{update\_example1}} : \coqref{SECF.Maps.examplemap'}{\coqdocdefinition{examplemap'}} "baz" \coqexternalref{::type scope:x '=' x}{http://rocq-prover.org/doc/V9.1.1/stdlib//Corelib.Init.Logic}{\coqdocnotation{=}} \coqexternalref{false}{http://rocq-prover.org/doc/V9.1.1/stdlib//Corelib.Init.Datatypes}{\coqdocconstructor{false}}.\coqdoceol
\coqdocnoindent
\coqdockw{Proof}. \coqdoctac{reflexivity}. \coqdockw{Qed}.\coqdoceol
\coqdocemptyline
\coqdocnoindent
\coqdockw{Example} \coqdef{SECF.Maps.update example2}{update\_example2}{\coqdocdefinition{update\_example2}} : \coqref{SECF.Maps.examplemap'}{\coqdocdefinition{examplemap'}} "foo" \coqexternalref{::type scope:x '=' x}{http://rocq-prover.org/doc/V9.1.1/stdlib//Corelib.Init.Logic}{\coqdocnotation{=}} \coqexternalref{true}{http://rocq-prover.org/doc/V9.1.1/stdlib//Corelib.Init.Datatypes}{\coqdocconstructor{true}}.\coqdoceol
\coqdocnoindent
\coqdockw{Proof}. \coqdoctac{reflexivity}. \coqdockw{Qed}.\coqdoceol
\coqdocemptyline
\coqdocnoindent
\coqdockw{Example} \coqdef{SECF.Maps.update example3}{update\_example3}{\coqdocdefinition{update\_example3}} : \coqref{SECF.Maps.examplemap'}{\coqdocdefinition{examplemap'}} "quux" \coqexternalref{::type scope:x '=' x}{http://rocq-prover.org/doc/V9.1.1/stdlib//Corelib.Init.Logic}{\coqdocnotation{=}} \coqexternalref{false}{http://rocq-prover.org/doc/V9.1.1/stdlib//Corelib.Init.Datatypes}{\coqdocconstructor{false}}.\coqdoceol
\coqdocnoindent
\coqdockw{Proof}. \coqdoctac{reflexivity}. \coqdockw{Qed}.\coqdoceol
\coqdocemptyline
\coqdocnoindent
\coqdockw{Example} \coqdef{SECF.Maps.update example4}{update\_example4}{\coqdocdefinition{update\_example4}} : \coqref{SECF.Maps.examplemap'}{\coqdocdefinition{examplemap'}} "bar" \coqexternalref{::type scope:x '=' x}{http://rocq-prover.org/doc/V9.1.1/stdlib//Corelib.Init.Logic}{\coqdocnotation{=}} \coqexternalref{true}{http://rocq-prover.org/doc/V9.1.1/stdlib//Corelib.Init.Datatypes}{\coqdocconstructor{true}}.\coqdoceol
\coqdocnoindent
\coqdockw{Proof}. \coqdoctac{reflexivity}. \coqdockw{Qed}.\coqdoceol
\coqdocemptyline
\end{coqdoccode}
When we use maps in later chapters, we'll need several fundamental
    facts about how they behave. 

 Even if you don't work the following exercises, make sure
    you thoroughly understand the statements of the lemmas! 

 (Some of the proofs require the functional extensionality axiom,
    which was discussed in the \coqdoclibrary{Logic} chapter.) 

\paragraph{Exercise: 1 star, standard, optional (t\_apply\_empty)}




    First, the empty map returns its default element for all keys: 
\begin{coqdoccode}
\coqdocemptyline
\coqdocnoindent
\coqdockw{Lemma} \coqdef{SECF.Maps.t apply empty}{t\_apply\_empty}{\coqdoclemma{t\_apply\_empty}} : \coqdockw{\ensuremath{\forall}} (\coqdef{SECF.Maps.A:16}{A}{\coqdocbinder{A}} : \coqdockw{Type}) (\coqdef{SECF.Maps.x:17}{x}{\coqdocbinder{x}} : \coqdocinductive{string}) (\coqdef{SECF.Maps.v:18}{v}{\coqdocbinder{v}} : \coqref{SECF.Maps.A:16}{\coqdocvariable{A}}),\coqdoceol
\coqdocindent{1.00em}
\coqref{SECF.Maps.:::'  ' '!->' x}{\coqdocnotation{(}}\coqref{SECF.Maps.:::'  ' '!->' x}{\coqdocnotation{\_\_}} \coqref{SECF.Maps.:::'  ' '!->' x}{\coqdocnotation{!->}} \coqref{SECF.Maps.v:18}{\coqdocvariable{v}}\coqref{SECF.Maps.:::'  ' '!->' x}{\coqdocnotation{)}} \coqref{SECF.Maps.:::'  ' '!->' x}{\coqdocnotation{x}} \coqexternalref{::type scope:x '=' x}{http://rocq-prover.org/doc/V9.1.1/stdlib//Corelib.Init.Logic}{\coqdocnotation{=}} \coqref{SECF.Maps.v:18}{\coqdocvariable{v}}.\coqdoceol
\coqdocnoindent
\coqdockw{Proof}.\coqdoceol
 \coqdocvar{Admitted}.\coqdoceol
\end{coqdoccode}
\ensuremath{\Box} 

\paragraph{Exercise: 2 stars, standard, optional (t\_update\_eq)}




    Next, if we update a map \coqdocvar{m} at a key \coqdocvar{x} with a new value \coqdocvar{v}
    and then look up \coqdocvar{x} in the map resulting from the \coqref{SECF.Maps.update}{\coqdocdefinition{update}}, we
    get back \coqdocvar{v}: 
\begin{coqdoccode}
\coqdocemptyline
\coqdocnoindent
\coqdockw{Lemma} \coqdef{SECF.Maps.t update eq}{t\_update\_eq}{\coqdoclemma{t\_update\_eq}} : \coqdockw{\ensuremath{\forall}} (\coqdef{SECF.Maps.A:19}{A}{\coqdocbinder{A}} : \coqdockw{Type}) (\coqdef{SECF.Maps.m:20}{m}{\coqdocbinder{m}} : \coqref{SECF.Maps.total map}{\coqdocdefinition{total\_map}} \coqref{SECF.Maps.A:19}{\coqdocvariable{A}}) \coqdef{SECF.Maps.x:21}{x}{\coqdocbinder{x}} \coqdef{SECF.Maps.v:22}{v}{\coqdocbinder{v}},\coqdoceol
\coqdocindent{1.00em}
\coqref{SECF.Maps.:::x '!->' x ';' x}{\coqdocnotation{(}}\coqref{SECF.Maps.x:21}{\coqdocvariable{x}} \coqref{SECF.Maps.:::x '!->' x ';' x}{\coqdocnotation{!->}} \coqref{SECF.Maps.v:22}{\coqdocvariable{v}} \coqref{SECF.Maps.:::x '!->' x ';' x}{\coqdocnotation{;}} \coqref{SECF.Maps.m:20}{\coqdocvariable{m}}\coqref{SECF.Maps.:::x '!->' x ';' x}{\coqdocnotation{)}} \coqref{SECF.Maps.:::x '!->' x ';' x}{\coqdocnotation{x}} \coqexternalref{::type scope:x '=' x}{http://rocq-prover.org/doc/V9.1.1/stdlib//Corelib.Init.Logic}{\coqdocnotation{=}} \coqref{SECF.Maps.v:22}{\coqdocvariable{v}}.\coqdoceol
\coqdocnoindent
\coqdockw{Proof}.\coqdoceol
 \coqdocvar{Admitted}.\coqdoceol
\end{coqdoccode}
\ensuremath{\Box} 

\paragraph{Exercise: 2 stars, standard, optional (t\_update\_neq)}




    On the other hand, if we update a map \coqdocvar{m} at a key \coqdocvar{x1} and then
    look up a \textit{different} key \coqdocvar{x2} in the resulting map, we get the
    same result that \coqdocvar{m} would have given: 
\begin{coqdoccode}
\coqdocemptyline
\coqdocnoindent
\coqdockw{Theorem} \coqdef{SECF.Maps.t update neq}{t\_update\_neq}{\coqdoclemma{t\_update\_neq}} : \coqdockw{\ensuremath{\forall}} (\coqdef{SECF.Maps.A:23}{A}{\coqdocbinder{A}} : \coqdockw{Type}) (\coqdef{SECF.Maps.m:24}{m}{\coqdocbinder{m}} : \coqref{SECF.Maps.total map}{\coqdocdefinition{total\_map}} \coqref{SECF.Maps.A:23}{\coqdocvariable{A}}) \coqdef{SECF.Maps.x1:25}{x1}{\coqdocbinder{x1}} \coqdef{SECF.Maps.x2:26}{x2}{\coqdocbinder{x2}} \coqdef{SECF.Maps.v:27}{v}{\coqdocbinder{v}},\coqdoceol
\coqdocindent{1.00em}
\coqref{SECF.Maps.x1:25}{\coqdocvariable{x1}} \coqexternalref{::type scope:x '<>' x}{http://rocq-prover.org/doc/V9.1.1/stdlib//Corelib.Init.Logic}{\coqdocnotation{\ensuremath{\not=}}} \coqref{SECF.Maps.x2:26}{\coqdocvariable{x2}} \coqexternalref{::type scope:x '->' x}{http://rocq-prover.org/doc/V9.1.1/stdlib//Corelib.Init.Logic}{\coqdocnotation{\ensuremath{\rightarrow}}}\coqdoceol
\coqdocindent{1.00em}
\coqref{SECF.Maps.:::x '!->' x ';' x}{\coqdocnotation{(}}\coqref{SECF.Maps.x1:25}{\coqdocvariable{x1}} \coqref{SECF.Maps.:::x '!->' x ';' x}{\coqdocnotation{!->}} \coqref{SECF.Maps.v:27}{\coqdocvariable{v}} \coqref{SECF.Maps.:::x '!->' x ';' x}{\coqdocnotation{;}} \coqref{SECF.Maps.m:24}{\coqdocvariable{m}}\coqref{SECF.Maps.:::x '!->' x ';' x}{\coqdocnotation{)}} \coqref{SECF.Maps.:::x '!->' x ';' x}{\coqdocnotation{x2}} \coqexternalref{::type scope:x '=' x}{http://rocq-prover.org/doc/V9.1.1/stdlib//Corelib.Init.Logic}{\coqdocnotation{=}} \coqref{SECF.Maps.m:24}{\coqdocvariable{m}} \coqref{SECF.Maps.x2:26}{\coqdocvariable{x2}}.\coqdoceol
\coqdocnoindent
\coqdockw{Proof}.\coqdoceol
 \coqdocvar{Admitted}.\coqdoceol
\end{coqdoccode}
\ensuremath{\Box} 

\paragraph{Exercise: 2 stars, standard, optional (t\_update\_shadow)}




    If we update a map \coqdocvar{m} at a key \coqdocvar{x} with a value \coqdocvar{v1} and then
    update again with the same key \coqdocvar{x} and another value \coqdocvar{v2}, the
    resulting map behaves the same (gives the same result when applied
    to any key) as the simpler map obtained by performing just
    the second \coqref{SECF.Maps.update}{\coqdocdefinition{update}} on \coqdocvar{m}: 
\begin{coqdoccode}
\coqdocemptyline
\coqdocnoindent
\coqdockw{Lemma} \coqdef{SECF.Maps.t update shadow}{t\_update\_shadow}{\coqdoclemma{t\_update\_shadow}} : \coqdockw{\ensuremath{\forall}} (\coqdef{SECF.Maps.A:28}{A}{\coqdocbinder{A}} : \coqdockw{Type}) (\coqdef{SECF.Maps.m:29}{m}{\coqdocbinder{m}} : \coqref{SECF.Maps.total map}{\coqdocdefinition{total\_map}} \coqref{SECF.Maps.A:28}{\coqdocvariable{A}}) \coqdef{SECF.Maps.x:30}{x}{\coqdocbinder{x}} \coqdef{SECF.Maps.v1:31}{v1}{\coqdocbinder{v1}} \coqdef{SECF.Maps.v2:32}{v2}{\coqdocbinder{v2}},\coqdoceol
\coqdocindent{1.00em}
\coqexternalref{::type scope:x '=' x}{http://rocq-prover.org/doc/V9.1.1/stdlib//Corelib.Init.Logic}{\coqdocnotation{(}}\coqref{SECF.Maps.x:30}{\coqdocvariable{x}} \coqref{SECF.Maps.:::x '!->' x ';' x}{\coqdocnotation{!->}} \coqref{SECF.Maps.v2:32}{\coqdocvariable{v2}} \coqref{SECF.Maps.:::x '!->' x ';' x}{\coqdocnotation{;}} \coqref{SECF.Maps.x:30}{\coqdocvariable{x}} \coqref{SECF.Maps.:::x '!->' x ';' x}{\coqdocnotation{!->}} \coqref{SECF.Maps.v1:31}{\coqdocvariable{v1}} \coqref{SECF.Maps.:::x '!->' x ';' x}{\coqdocnotation{;}} \coqref{SECF.Maps.m:29}{\coqdocvariable{m}}\coqexternalref{::type scope:x '=' x}{http://rocq-prover.org/doc/V9.1.1/stdlib//Corelib.Init.Logic}{\coqdocnotation{)}} \coqexternalref{::type scope:x '=' x}{http://rocq-prover.org/doc/V9.1.1/stdlib//Corelib.Init.Logic}{\coqdocnotation{=}} \coqexternalref{::type scope:x '=' x}{http://rocq-prover.org/doc/V9.1.1/stdlib//Corelib.Init.Logic}{\coqdocnotation{(}}\coqref{SECF.Maps.x:30}{\coqdocvariable{x}} \coqref{SECF.Maps.:::x '!->' x ';' x}{\coqdocnotation{!->}} \coqref{SECF.Maps.v2:32}{\coqdocvariable{v2}} \coqref{SECF.Maps.:::x '!->' x ';' x}{\coqdocnotation{;}} \coqref{SECF.Maps.m:29}{\coqdocvariable{m}}\coqexternalref{::type scope:x '=' x}{http://rocq-prover.org/doc/V9.1.1/stdlib//Corelib.Init.Logic}{\coqdocnotation{)}}.\coqdoceol
\coqdocnoindent
\coqdockw{Proof}.\coqdoceol
 \coqdocvar{Admitted}.\coqdoceol
\end{coqdoccode}
\ensuremath{\Box} 

\paragraph{Exercise: 2 stars, standard (t\_update\_same)}




    Given \coqdocinductive{string}s \coqdocvar{x1} and \coqdocvar{x2}, we can use the tactic
    \coqdoctac{destruct} (\coqdoclemma{eqb\_spec} \coqdocvar{x1} \coqdocvar{x2}) to simultaneously perform case
    analysis on the result of \coqdocvar{String.eqb} \coqdocvar{x1} \coqdocvar{x2} and generate
    hypotheses about the equality (in the sense of =) of \coqdocvar{x1} and
    \coqdocvar{x2}.  With the example in chapter \coqdocvar{IndProp} as a template,
    use \coqdocvar{String.eqb\_spec} to prove the following theorem, which states
    that if we update a map to assign key \coqdocvar{x} the same value as it
    already has in \coqdocvar{m}, then the result is equal to \coqdocvar{m}: 
\begin{coqdoccode}
\coqdocemptyline
\coqdocnoindent
\coqdockw{Theorem} \coqdef{SECF.Maps.t update same}{t\_update\_same}{\coqdoclemma{t\_update\_same}} : \coqdockw{\ensuremath{\forall}} (\coqdef{SECF.Maps.A:33}{A}{\coqdocbinder{A}} : \coqdockw{Type}) (\coqdef{SECF.Maps.m:34}{m}{\coqdocbinder{m}} : \coqref{SECF.Maps.total map}{\coqdocdefinition{total\_map}} \coqref{SECF.Maps.A:33}{\coqdocvariable{A}}) \coqdef{SECF.Maps.x:35}{x}{\coqdocbinder{x}},\coqdoceol
\coqdocindent{1.00em}
\coqexternalref{::type scope:x '=' x}{http://rocq-prover.org/doc/V9.1.1/stdlib//Corelib.Init.Logic}{\coqdocnotation{(}}\coqref{SECF.Maps.x:35}{\coqdocvariable{x}} \coqref{SECF.Maps.:::x '!->' x ';' x}{\coqdocnotation{!->}} \coqref{SECF.Maps.m:34}{\coqdocvariable{m}} \coqref{SECF.Maps.x:35}{\coqdocvariable{x}} \coqref{SECF.Maps.:::x '!->' x ';' x}{\coqdocnotation{;}} \coqref{SECF.Maps.m:34}{\coqdocvariable{m}}\coqexternalref{::type scope:x '=' x}{http://rocq-prover.org/doc/V9.1.1/stdlib//Corelib.Init.Logic}{\coqdocnotation{)}} \coqexternalref{::type scope:x '=' x}{http://rocq-prover.org/doc/V9.1.1/stdlib//Corelib.Init.Logic}{\coqdocnotation{=}} \coqref{SECF.Maps.m:34}{\coqdocvariable{m}}.\coqdoceol
\coqdocnoindent
\coqdockw{Proof}.\coqdoceol
 \coqdocvar{Admitted}.\coqdoceol
\end{coqdoccode}
\ensuremath{\Box} 

\paragraph{Exercise: 3 stars, standard, especially useful (t\_update\_permute)}




    Similarly, use \coqdocvar{String.eqb\_spec} to prove one final property of
    the \coqref{SECF.Maps.update}{\coqdocdefinition{update}} function: If we update a map \coqdocvar{m} at two distinct
    keys, it doesn't matter in which order we do the updates. 
\begin{coqdoccode}
\coqdocemptyline
\coqdocnoindent
\coqdockw{Theorem} \coqdef{SECF.Maps.t update permute}{t\_update\_permute}{\coqdoclemma{t\_update\_permute}} : \coqdockw{\ensuremath{\forall}} (\coqdef{SECF.Maps.A:36}{A}{\coqdocbinder{A}} : \coqdockw{Type}) (\coqdef{SECF.Maps.m:37}{m}{\coqdocbinder{m}} : \coqref{SECF.Maps.total map}{\coqdocdefinition{total\_map}} \coqref{SECF.Maps.A:36}{\coqdocvariable{A}})\coqdoceol
\coqdocindent{17.00em}
\coqdef{SECF.Maps.v1:38}{v1}{\coqdocbinder{v1}} \coqdef{SECF.Maps.v2:39}{v2}{\coqdocbinder{v2}} \coqdef{SECF.Maps.x1:40}{x1}{\coqdocbinder{x1}} \coqdef{SECF.Maps.x2:41}{x2}{\coqdocbinder{x2}},\coqdoceol
\coqdocindent{1.00em}
\coqref{SECF.Maps.x2:41}{\coqdocvariable{x2}} \coqexternalref{::type scope:x '<>' x}{http://rocq-prover.org/doc/V9.1.1/stdlib//Corelib.Init.Logic}{\coqdocnotation{\ensuremath{\not=}}} \coqref{SECF.Maps.x1:40}{\coqdocvariable{x1}} \coqexternalref{::type scope:x '->' x}{http://rocq-prover.org/doc/V9.1.1/stdlib//Corelib.Init.Logic}{\coqdocnotation{\ensuremath{\rightarrow}}}\coqdoceol
\coqdocindent{1.00em}
\coqexternalref{::type scope:x '=' x}{http://rocq-prover.org/doc/V9.1.1/stdlib//Corelib.Init.Logic}{\coqdocnotation{(}}\coqref{SECF.Maps.x1:40}{\coqdocvariable{x1}} \coqref{SECF.Maps.:::x '!->' x ';' x}{\coqdocnotation{!->}} \coqref{SECF.Maps.v1:38}{\coqdocvariable{v1}} \coqref{SECF.Maps.:::x '!->' x ';' x}{\coqdocnotation{;}} \coqref{SECF.Maps.x2:41}{\coqdocvariable{x2}} \coqref{SECF.Maps.:::x '!->' x ';' x}{\coqdocnotation{!->}} \coqref{SECF.Maps.v2:39}{\coqdocvariable{v2}} \coqref{SECF.Maps.:::x '!->' x ';' x}{\coqdocnotation{;}} \coqref{SECF.Maps.m:37}{\coqdocvariable{m}}\coqexternalref{::type scope:x '=' x}{http://rocq-prover.org/doc/V9.1.1/stdlib//Corelib.Init.Logic}{\coqdocnotation{)}}\coqdoceol
\coqdocindent{1.00em}
\coqexternalref{::type scope:x '=' x}{http://rocq-prover.org/doc/V9.1.1/stdlib//Corelib.Init.Logic}{\coqdocnotation{=}}\coqdoceol
\coqdocindent{1.00em}
\coqexternalref{::type scope:x '=' x}{http://rocq-prover.org/doc/V9.1.1/stdlib//Corelib.Init.Logic}{\coqdocnotation{(}}\coqref{SECF.Maps.x2:41}{\coqdocvariable{x2}} \coqref{SECF.Maps.:::x '!->' x ';' x}{\coqdocnotation{!->}} \coqref{SECF.Maps.v2:39}{\coqdocvariable{v2}} \coqref{SECF.Maps.:::x '!->' x ';' x}{\coqdocnotation{;}} \coqref{SECF.Maps.x1:40}{\coqdocvariable{x1}} \coqref{SECF.Maps.:::x '!->' x ';' x}{\coqdocnotation{!->}} \coqref{SECF.Maps.v1:38}{\coqdocvariable{v1}} \coqref{SECF.Maps.:::x '!->' x ';' x}{\coqdocnotation{;}} \coqref{SECF.Maps.m:37}{\coqdocvariable{m}}\coqexternalref{::type scope:x '=' x}{http://rocq-prover.org/doc/V9.1.1/stdlib//Corelib.Init.Logic}{\coqdocnotation{)}}.\coqdoceol
\coqdocnoindent
\coqdockw{Proof}.\coqdoceol
 \coqdocvar{Admitted}.\coqdoceol
\end{coqdoccode}
\ensuremath{\Box} 
\begin{coqdoccode}
\coqdocemptyline
\end{coqdoccode}
\section{Partial maps}



 Lastly, we define \textit{partial maps} on top of total maps.  A partial
    map with elements of type \coqdocvar{A} is simply a total map with elements
    of type \coqexternalref{option}{http://rocq-prover.org/doc/V9.1.1/stdlib//Corelib.Init.Datatypes}{\coqdocinductive{option}} \coqdocvar{A} and default element \coqexternalref{None}{http://rocq-prover.org/doc/V9.1.1/stdlib//Corelib.Init.Datatypes}{\coqdocconstructor{None}}. 
\begin{coqdoccode}
\coqdocemptyline
\coqdocnoindent
\coqdockw{Definition} \coqdef{SECF.Maps.partial map}{partial\_map}{\coqdocdefinition{partial\_map}} (\coqdef{SECF.Maps.A:42}{A}{\coqdocbinder{A}} : \coqdockw{Type}) := \coqref{SECF.Maps.total map}{\coqdocdefinition{total\_map}} (\coqexternalref{option}{http://rocq-prover.org/doc/V9.1.1/stdlib//Corelib.Init.Datatypes}{\coqdocinductive{option}} \coqref{SECF.Maps.A:42}{\coqdocvariable{A}}).\coqdoceol
\coqdocemptyline
\coqdocnoindent
\coqdockw{Definition} \coqdef{SECF.Maps.empty}{empty}{\coqdocdefinition{empty}} \{\coqdef{SECF.Maps.A:43}{A}{\coqdocbinder{A}} : \coqdockw{Type}\} : \coqref{SECF.Maps.partial map}{\coqdocdefinition{partial\_map}} \coqref{SECF.Maps.A:43}{\coqdocvariable{A}} :=\coqdoceol
\coqdocindent{1.00em}
\coqref{SECF.Maps.t empty}{\coqdocdefinition{t\_empty}} \coqexternalref{None}{http://rocq-prover.org/doc/V9.1.1/stdlib//Corelib.Init.Datatypes}{\coqdocconstructor{None}}.\coqdoceol
\coqdocemptyline
\coqdocnoindent
\coqdockw{Definition} \coqdef{SECF.Maps.update}{update}{\coqdocdefinition{update}} \{\coqdef{SECF.Maps.A:44}{A}{\coqdocbinder{A}} : \coqdockw{Type}\} (\coqdef{SECF.Maps.m:45}{m}{\coqdocbinder{m}} : \coqref{SECF.Maps.partial map}{\coqdocdefinition{partial\_map}} \coqref{SECF.Maps.A:44}{\coqdocvariable{A}})\coqdoceol
\coqdocindent{5.50em}
(\coqdef{SECF.Maps.x:46}{x}{\coqdocbinder{x}} : \coqdocinductive{string}) (\coqdef{SECF.Maps.v:47}{v}{\coqdocbinder{v}} : \coqref{SECF.Maps.A:44}{\coqdocvariable{A}}) :=\coqdoceol
\coqdocindent{1.00em}
(\coqref{SECF.Maps.x:46}{\coqdocvariable{x}} \coqref{SECF.Maps.:::x '!->' x ';' x}{\coqdocnotation{!->}} \coqexternalref{Some}{http://rocq-prover.org/doc/V9.1.1/stdlib//Corelib.Init.Datatypes}{\coqdocconstructor{Some}} \coqref{SECF.Maps.v:47}{\coqdocvariable{v}} \coqref{SECF.Maps.:::x '!->' x ';' x}{\coqdocnotation{;}} \coqref{SECF.Maps.m:45}{\coqdocvariable{m}}).\coqdoceol
\end{coqdoccode}
We introduce a similar notation for partial maps: 
\begin{coqdoccode}
\coqdocnoindent
\coqdockw{Notation} \coqdef{SECF.Maps.:::x '|->' x ';' x}{"}{"}x '|->' v ';' m" := (\coqref{SECF.Maps.update}{\coqdocdefinition{update}} \coqdocvar{m} \coqdocvar{x} \coqdocvar{v})\coqdoceol
\coqdocindent{1.00em}
(\coqdoctac{at} \coqdockw{level} 0, \coqdocvar{x} \coqdockw{constr}, \coqdocvar{v} \coqdoctac{at} \coqdockw{level} 200, \coqdoctac{right} \coqdockw{associativity}).\coqdoceol
\coqdocemptyline
\end{coqdoccode}
We can also hide the last case when it is empty. 
\begin{coqdoccode}
\coqdocnoindent
\coqdockw{Notation} \coqdef{SECF.Maps.:::x '|->' x}{"}{"}x '|->' v" := (\coqref{SECF.Maps.update}{\coqdocdefinition{update}} \coqref{SECF.Maps.empty}{\coqdocdefinition{empty}} \coqdocvar{x} \coqdocvar{v})\coqdoceol
\coqdocindent{1.00em}
(\coqdoctac{at} \coqdockw{level} 0, \coqdocvar{x} \coqdockw{constr}, \coqdocvar{v} \coqdoctac{at} \coqdockw{level} 200).\coqdoceol
\coqdocemptyline
\coqdocnoindent
\coqdockw{Definition} \coqdef{SECF.Maps.examplepmap}{examplepmap}{\coqdocdefinition{examplepmap}} :=\coqdoceol
\coqdocindent{1.00em}
("Church" \coqref{SECF.Maps.:::x '|->' x ';' x}{\coqdocnotation{|->}} \coqexternalref{true}{http://rocq-prover.org/doc/V9.1.1/stdlib//Corelib.Init.Datatypes}{\coqdocconstructor{true}} \coqref{SECF.Maps.:::x '|->' x ';' x}{\coqdocnotation{;}} "Turing" \coqref{SECF.Maps.:::x '|->' x}{\coqdocnotation{|->}} \coqexternalref{false}{http://rocq-prover.org/doc/V9.1.1/stdlib//Corelib.Init.Datatypes}{\coqdocconstructor{false}}).\coqdoceol
\coqdocemptyline
\end{coqdoccode}
We now straightforwardly lift all of the basic lemmas about total
    maps to partial maps.  
\begin{coqdoccode}
\coqdocemptyline
\coqdocnoindent
\coqdockw{Lemma} \coqdef{SECF.Maps.apply empty}{apply\_empty}{\coqdoclemma{apply\_empty}} : \coqdockw{\ensuremath{\forall}} (\coqdef{SECF.Maps.A:48}{A}{\coqdocbinder{A}} : \coqdockw{Type}) (\coqdef{SECF.Maps.x:49}{x}{\coqdocbinder{x}} : \coqdocinductive{string}),\coqdoceol
\coqdocindent{1.00em}
@\coqref{SECF.Maps.empty}{\coqdocdefinition{empty}} \coqref{SECF.Maps.A:48}{\coqdocvariable{A}} \coqref{SECF.Maps.x:49}{\coqdocvariable{x}} \coqexternalref{::type scope:x '=' x}{http://rocq-prover.org/doc/V9.1.1/stdlib//Corelib.Init.Logic}{\coqdocnotation{=}} \coqexternalref{None}{http://rocq-prover.org/doc/V9.1.1/stdlib//Corelib.Init.Datatypes}{\coqdocconstructor{None}}.\coqdoceol
\coqdocnoindent
\coqdockw{Proof}.\coqdoceol
\coqdocindent{1.00em}
\coqdoctac{intros}. \coqdoctac{unfold} \coqref{SECF.Maps.empty}{\coqdocdefinition{empty}}. \coqdoctac{rewrite} \coqref{SECF.Maps.t apply empty}{\coqdocaxiom{t\_apply\_empty}}.\coqdoceol
\coqdocindent{1.00em}
\coqdoctac{reflexivity}.\coqdoceol
\coqdocnoindent
\coqdockw{Qed}.\coqdoceol
\coqdocemptyline
\coqdocnoindent
\coqdockw{Lemma} \coqdef{SECF.Maps.update eq}{update\_eq}{\coqdoclemma{update\_eq}} : \coqdockw{\ensuremath{\forall}} (\coqdef{SECF.Maps.A:50}{A}{\coqdocbinder{A}} : \coqdockw{Type}) (\coqdef{SECF.Maps.m:51}{m}{\coqdocbinder{m}} : \coqref{SECF.Maps.partial map}{\coqdocdefinition{partial\_map}} \coqref{SECF.Maps.A:50}{\coqdocvariable{A}}) \coqdef{SECF.Maps.x:52}{x}{\coqdocbinder{x}} \coqdef{SECF.Maps.v:53}{v}{\coqdocbinder{v}},\coqdoceol
\coqdocindent{1.00em}
\coqref{SECF.Maps.:::x '|->' x ';' x}{\coqdocnotation{(}}\coqref{SECF.Maps.x:52}{\coqdocvariable{x}} \coqref{SECF.Maps.:::x '|->' x ';' x}{\coqdocnotation{|->}} \coqref{SECF.Maps.v:53}{\coqdocvariable{v}} \coqref{SECF.Maps.:::x '|->' x ';' x}{\coqdocnotation{;}} \coqref{SECF.Maps.m:51}{\coqdocvariable{m}}\coqref{SECF.Maps.:::x '|->' x ';' x}{\coqdocnotation{)}} \coqref{SECF.Maps.:::x '|->' x ';' x}{\coqdocnotation{x}} \coqexternalref{::type scope:x '=' x}{http://rocq-prover.org/doc/V9.1.1/stdlib//Corelib.Init.Logic}{\coqdocnotation{=}} \coqexternalref{Some}{http://rocq-prover.org/doc/V9.1.1/stdlib//Corelib.Init.Datatypes}{\coqdocconstructor{Some}} \coqref{SECF.Maps.v:53}{\coqdocvariable{v}}.\coqdoceol
\coqdocnoindent
\coqdockw{Proof}.\coqdoceol
\coqdocindent{1.00em}
\coqdoctac{intros}. \coqdoctac{unfold} \coqref{SECF.Maps.update}{\coqdocdefinition{update}}. \coqdoctac{rewrite} \coqref{SECF.Maps.t update eq}{\coqdocaxiom{t\_update\_eq}}.\coqdoceol
\coqdocindent{1.00em}
\coqdoctac{reflexivity}.\coqdoceol
\coqdocnoindent
\coqdockw{Qed}.\coqdoceol
\coqdocemptyline
\end{coqdoccode}
The \coqref{SECF.Maps.update eq}{\coqdoclemma{update\_eq}} lemma is used very often in proofs.  Adding it to
    Rocq's global ``hint database'' allows proof-automation tactics such
    as \coqdoctac{auto} to find it. 
\begin{coqdoccode}
\coqdocnoindent
\#[\coqdocvar{global}] \coqdockw{Hint} \coqdockw{Resolve} \coqdocvar{update\_eq} : \coqdocvar{core}.\coqdoceol
\coqdocemptyline
\coqdocnoindent
\coqdockw{Theorem} \coqdef{SECF.Maps.update neq}{update\_neq}{\coqdoclemma{update\_neq}} : \coqdockw{\ensuremath{\forall}} (\coqdef{SECF.Maps.A:54}{A}{\coqdocbinder{A}} : \coqdockw{Type}) (\coqdef{SECF.Maps.m:55}{m}{\coqdocbinder{m}} : \coqref{SECF.Maps.partial map}{\coqdocdefinition{partial\_map}} \coqref{SECF.Maps.A:54}{\coqdocvariable{A}}) \coqdef{SECF.Maps.x1:56}{x1}{\coqdocbinder{x1}} \coqdef{SECF.Maps.x2:57}{x2}{\coqdocbinder{x2}} \coqdef{SECF.Maps.v:58}{v}{\coqdocbinder{v}},\coqdoceol
\coqdocindent{1.00em}
\coqref{SECF.Maps.x2:57}{\coqdocvariable{x2}} \coqexternalref{::type scope:x '<>' x}{http://rocq-prover.org/doc/V9.1.1/stdlib//Corelib.Init.Logic}{\coqdocnotation{\ensuremath{\not=}}} \coqref{SECF.Maps.x1:56}{\coqdocvariable{x1}} \coqexternalref{::type scope:x '->' x}{http://rocq-prover.org/doc/V9.1.1/stdlib//Corelib.Init.Logic}{\coqdocnotation{\ensuremath{\rightarrow}}}\coqdoceol
\coqdocindent{1.00em}
\coqref{SECF.Maps.:::x '|->' x ';' x}{\coqdocnotation{(}}\coqref{SECF.Maps.x2:57}{\coqdocvariable{x2}} \coqref{SECF.Maps.:::x '|->' x ';' x}{\coqdocnotation{|->}} \coqref{SECF.Maps.v:58}{\coqdocvariable{v}} \coqref{SECF.Maps.:::x '|->' x ';' x}{\coqdocnotation{;}} \coqref{SECF.Maps.m:55}{\coqdocvariable{m}}\coqref{SECF.Maps.:::x '|->' x ';' x}{\coqdocnotation{)}} \coqref{SECF.Maps.:::x '|->' x ';' x}{\coqdocnotation{x1}} \coqexternalref{::type scope:x '=' x}{http://rocq-prover.org/doc/V9.1.1/stdlib//Corelib.Init.Logic}{\coqdocnotation{=}} \coqref{SECF.Maps.m:55}{\coqdocvariable{m}} \coqref{SECF.Maps.x1:56}{\coqdocvariable{x1}}.\coqdoceol
\coqdocnoindent
\coqdockw{Proof}.\coqdoceol
\coqdocindent{1.00em}
\coqdoctac{intros} \coqdocvar{A} \coqdocvar{m} \coqdocvar{x1} \coqdocvar{x2} \coqdocvar{v} \coqdocvar{H}.\coqdoceol
\coqdocindent{1.00em}
\coqdoctac{unfold} \coqref{SECF.Maps.update}{\coqdocdefinition{update}}. \coqdoctac{rewrite} \coqref{SECF.Maps.t update neq}{\coqdocaxiom{t\_update\_neq}}.\coqdoceol
\coqdocindent{1.00em}
- \coqdoctac{reflexivity}.\coqdoceol
\coqdocindent{1.00em}
- \coqdoctac{apply} \coqdocvar{H}.\coqdoceol
\coqdocnoindent
\coqdockw{Qed}.\coqdoceol
\coqdocemptyline
\coqdocnoindent
\coqdockw{Lemma} \coqdef{SECF.Maps.update shadow}{update\_shadow}{\coqdoclemma{update\_shadow}} : \coqdockw{\ensuremath{\forall}} (\coqdef{SECF.Maps.A:59}{A}{\coqdocbinder{A}} : \coqdockw{Type}) (\coqdef{SECF.Maps.m:60}{m}{\coqdocbinder{m}} : \coqref{SECF.Maps.partial map}{\coqdocdefinition{partial\_map}} \coqref{SECF.Maps.A:59}{\coqdocvariable{A}}) \coqdef{SECF.Maps.x:61}{x}{\coqdocbinder{x}} \coqdef{SECF.Maps.v1:62}{v1}{\coqdocbinder{v1}} \coqdef{SECF.Maps.v2:63}{v2}{\coqdocbinder{v2}},\coqdoceol
\coqdocindent{1.00em}
\coqexternalref{::type scope:x '=' x}{http://rocq-prover.org/doc/V9.1.1/stdlib//Corelib.Init.Logic}{\coqdocnotation{(}}\coqref{SECF.Maps.x:61}{\coqdocvariable{x}} \coqref{SECF.Maps.:::x '|->' x ';' x}{\coqdocnotation{|->}} \coqref{SECF.Maps.v2:63}{\coqdocvariable{v2}} \coqref{SECF.Maps.:::x '|->' x ';' x}{\coqdocnotation{;}} \coqref{SECF.Maps.x:61}{\coqdocvariable{x}} \coqref{SECF.Maps.:::x '|->' x ';' x}{\coqdocnotation{|->}} \coqref{SECF.Maps.v1:62}{\coqdocvariable{v1}} \coqref{SECF.Maps.:::x '|->' x ';' x}{\coqdocnotation{;}} \coqref{SECF.Maps.m:60}{\coqdocvariable{m}}\coqexternalref{::type scope:x '=' x}{http://rocq-prover.org/doc/V9.1.1/stdlib//Corelib.Init.Logic}{\coqdocnotation{)}} \coqexternalref{::type scope:x '=' x}{http://rocq-prover.org/doc/V9.1.1/stdlib//Corelib.Init.Logic}{\coqdocnotation{=}} \coqexternalref{::type scope:x '=' x}{http://rocq-prover.org/doc/V9.1.1/stdlib//Corelib.Init.Logic}{\coqdocnotation{(}}\coqref{SECF.Maps.x:61}{\coqdocvariable{x}} \coqref{SECF.Maps.:::x '|->' x ';' x}{\coqdocnotation{|->}} \coqref{SECF.Maps.v2:63}{\coqdocvariable{v2}} \coqref{SECF.Maps.:::x '|->' x ';' x}{\coqdocnotation{;}} \coqref{SECF.Maps.m:60}{\coqdocvariable{m}}\coqexternalref{::type scope:x '=' x}{http://rocq-prover.org/doc/V9.1.1/stdlib//Corelib.Init.Logic}{\coqdocnotation{)}}.\coqdoceol
\coqdocnoindent
\coqdockw{Proof}.\coqdoceol
\coqdocindent{1.00em}
\coqdoctac{intros} \coqdocvar{A} \coqdocvar{m} \coqdocvar{x} \coqdocvar{v1} \coqdocvar{v2}. \coqdoctac{unfold} \coqref{SECF.Maps.update}{\coqdocdefinition{update}}. \coqdoctac{rewrite} \coqref{SECF.Maps.t update shadow}{\coqdocaxiom{t\_update\_shadow}}.\coqdoceol
\coqdocindent{1.00em}
\coqdoctac{reflexivity}.\coqdoceol
\coqdocnoindent
\coqdockw{Qed}.\coqdoceol
\coqdocemptyline
\coqdocnoindent
\coqdockw{Theorem} \coqdef{SECF.Maps.update same}{update\_same}{\coqdoclemma{update\_same}} : \coqdockw{\ensuremath{\forall}} (\coqdef{SECF.Maps.A:64}{A}{\coqdocbinder{A}} : \coqdockw{Type}) (\coqdef{SECF.Maps.m:65}{m}{\coqdocbinder{m}} : \coqref{SECF.Maps.partial map}{\coqdocdefinition{partial\_map}} \coqref{SECF.Maps.A:64}{\coqdocvariable{A}}) \coqdef{SECF.Maps.x:66}{x}{\coqdocbinder{x}} \coqdef{SECF.Maps.v:67}{v}{\coqdocbinder{v}},\coqdoceol
\coqdocindent{1.00em}
\coqref{SECF.Maps.m:65}{\coqdocvariable{m}} \coqref{SECF.Maps.x:66}{\coqdocvariable{x}} \coqexternalref{::type scope:x '=' x}{http://rocq-prover.org/doc/V9.1.1/stdlib//Corelib.Init.Logic}{\coqdocnotation{=}} \coqexternalref{Some}{http://rocq-prover.org/doc/V9.1.1/stdlib//Corelib.Init.Datatypes}{\coqdocconstructor{Some}} \coqref{SECF.Maps.v:67}{\coqdocvariable{v}} \coqexternalref{::type scope:x '->' x}{http://rocq-prover.org/doc/V9.1.1/stdlib//Corelib.Init.Logic}{\coqdocnotation{\ensuremath{\rightarrow}}}\coqdoceol
\coqdocindent{1.00em}
\coqexternalref{::type scope:x '=' x}{http://rocq-prover.org/doc/V9.1.1/stdlib//Corelib.Init.Logic}{\coqdocnotation{(}}\coqref{SECF.Maps.x:66}{\coqdocvariable{x}} \coqref{SECF.Maps.:::x '|->' x ';' x}{\coqdocnotation{|->}} \coqref{SECF.Maps.v:67}{\coqdocvariable{v}} \coqref{SECF.Maps.:::x '|->' x ';' x}{\coqdocnotation{;}} \coqref{SECF.Maps.m:65}{\coqdocvariable{m}}\coqexternalref{::type scope:x '=' x}{http://rocq-prover.org/doc/V9.1.1/stdlib//Corelib.Init.Logic}{\coqdocnotation{)}} \coqexternalref{::type scope:x '=' x}{http://rocq-prover.org/doc/V9.1.1/stdlib//Corelib.Init.Logic}{\coqdocnotation{=}} \coqref{SECF.Maps.m:65}{\coqdocvariable{m}}.\coqdoceol
\coqdocnoindent
\coqdockw{Proof}.\coqdoceol
\coqdocindent{1.00em}
\coqdoctac{intros} \coqdocvar{A} \coqdocvar{m} \coqdocvar{x} \coqdocvar{v} \coqdocvar{H}. \coqdoctac{unfold} \coqref{SECF.Maps.update}{\coqdocdefinition{update}}. \coqdoctac{rewrite} \ensuremath{\leftarrow} \coqdocvar{H}.\coqdoceol
\coqdocindent{1.00em}
\coqdoctac{apply} \coqref{SECF.Maps.t update same}{\coqdocaxiom{t\_update\_same}}.\coqdoceol
\coqdocnoindent
\coqdockw{Qed}.\coqdoceol
\coqdocemptyline
\coqdocnoindent
\coqdockw{Theorem} \coqdef{SECF.Maps.update permute}{update\_permute}{\coqdoclemma{update\_permute}} : \coqdockw{\ensuremath{\forall}} (\coqdef{SECF.Maps.A:68}{A}{\coqdocbinder{A}} : \coqdockw{Type}) (\coqdef{SECF.Maps.m:69}{m}{\coqdocbinder{m}} : \coqref{SECF.Maps.partial map}{\coqdocdefinition{partial\_map}} \coqref{SECF.Maps.A:68}{\coqdocvariable{A}})\coqdoceol
\coqdocindent{16.00em}
\coqdef{SECF.Maps.x1:70}{x1}{\coqdocbinder{x1}} \coqdef{SECF.Maps.x2:71}{x2}{\coqdocbinder{x2}} \coqdef{SECF.Maps.v1:72}{v1}{\coqdocbinder{v1}} \coqdef{SECF.Maps.v2:73}{v2}{\coqdocbinder{v2}},\coqdoceol
\coqdocindent{1.00em}
\coqref{SECF.Maps.x2:71}{\coqdocvariable{x2}} \coqexternalref{::type scope:x '<>' x}{http://rocq-prover.org/doc/V9.1.1/stdlib//Corelib.Init.Logic}{\coqdocnotation{\ensuremath{\not=}}} \coqref{SECF.Maps.x1:70}{\coqdocvariable{x1}} \coqexternalref{::type scope:x '->' x}{http://rocq-prover.org/doc/V9.1.1/stdlib//Corelib.Init.Logic}{\coqdocnotation{\ensuremath{\rightarrow}}}\coqdoceol
\coqdocindent{1.00em}
\coqexternalref{::type scope:x '=' x}{http://rocq-prover.org/doc/V9.1.1/stdlib//Corelib.Init.Logic}{\coqdocnotation{(}}\coqref{SECF.Maps.x1:70}{\coqdocvariable{x1}} \coqref{SECF.Maps.:::x '|->' x ';' x}{\coqdocnotation{|->}} \coqref{SECF.Maps.v1:72}{\coqdocvariable{v1}} \coqref{SECF.Maps.:::x '|->' x ';' x}{\coqdocnotation{;}} \coqref{SECF.Maps.x2:71}{\coqdocvariable{x2}} \coqref{SECF.Maps.:::x '|->' x ';' x}{\coqdocnotation{|->}} \coqref{SECF.Maps.v2:73}{\coqdocvariable{v2}} \coqref{SECF.Maps.:::x '|->' x ';' x}{\coqdocnotation{;}} \coqref{SECF.Maps.m:69}{\coqdocvariable{m}}\coqexternalref{::type scope:x '=' x}{http://rocq-prover.org/doc/V9.1.1/stdlib//Corelib.Init.Logic}{\coqdocnotation{)}} \coqexternalref{::type scope:x '=' x}{http://rocq-prover.org/doc/V9.1.1/stdlib//Corelib.Init.Logic}{\coqdocnotation{=}} \coqexternalref{::type scope:x '=' x}{http://rocq-prover.org/doc/V9.1.1/stdlib//Corelib.Init.Logic}{\coqdocnotation{(}}\coqref{SECF.Maps.x2:71}{\coqdocvariable{x2}} \coqref{SECF.Maps.:::x '|->' x ';' x}{\coqdocnotation{|->}} \coqref{SECF.Maps.v2:73}{\coqdocvariable{v2}} \coqref{SECF.Maps.:::x '|->' x ';' x}{\coqdocnotation{;}} \coqref{SECF.Maps.x1:70}{\coqdocvariable{x1}} \coqref{SECF.Maps.:::x '|->' x ';' x}{\coqdocnotation{|->}} \coqref{SECF.Maps.v1:72}{\coqdocvariable{v1}} \coqref{SECF.Maps.:::x '|->' x ';' x}{\coqdocnotation{;}} \coqref{SECF.Maps.m:69}{\coqdocvariable{m}}\coqexternalref{::type scope:x '=' x}{http://rocq-prover.org/doc/V9.1.1/stdlib//Corelib.Init.Logic}{\coqdocnotation{)}}.\coqdoceol
\coqdocnoindent
\coqdockw{Proof}.\coqdoceol
\coqdocindent{1.00em}
\coqdoctac{intros} \coqdocvar{A} \coqdocvar{m} \coqdocvar{x1} \coqdocvar{x2} \coqdocvar{v1} \coqdocvar{v2}. \coqdoctac{unfold} \coqref{SECF.Maps.update}{\coqdocdefinition{update}}.\coqdoceol
\coqdocindent{1.00em}
\coqdoctac{apply} \coqref{SECF.Maps.t update permute}{\coqdocaxiom{t\_update\_permute}}.\coqdoceol
\coqdocnoindent
\coqdockw{Qed}.\coqdoceol
\coqdocemptyline
\end{coqdoccode}
One last thing: For partial maps, it's convenient to introduce a
    notion of map inclusion, stating that all the entries in one map
    are also present in another: 
\begin{coqdoccode}
\coqdocemptyline
\coqdocnoindent
\coqdockw{Definition} \coqdef{SECF.Maps.includedin}{includedin}{\coqdocdefinition{includedin}} \{\coqdef{SECF.Maps.A:74}{A}{\coqdocbinder{A}} : \coqdockw{Type}\} (\coqdef{SECF.Maps.m:75}{m}{\coqdocbinder{m}} \coqdef{SECF.Maps.m':76}{m'}{\coqdocbinder{m'}} : \coqref{SECF.Maps.partial map}{\coqdocdefinition{partial\_map}} \coqref{SECF.Maps.A:74}{\coqdocvariable{A}}) :=\coqdoceol
\coqdocindent{1.00em}
\coqdockw{\ensuremath{\forall}} \coqdef{SECF.Maps.x:77}{x}{\coqdocbinder{x}} \coqdef{SECF.Maps.v:78}{v}{\coqdocbinder{v}}, \coqref{SECF.Maps.m:75}{\coqdocvariable{m}} \coqref{SECF.Maps.x:77}{\coqdocvariable{x}} \coqexternalref{::type scope:x '=' x}{http://rocq-prover.org/doc/V9.1.1/stdlib//Corelib.Init.Logic}{\coqdocnotation{=}} \coqexternalref{Some}{http://rocq-prover.org/doc/V9.1.1/stdlib//Corelib.Init.Datatypes}{\coqdocconstructor{Some}} \coqref{SECF.Maps.v:78}{\coqdocvariable{v}} \coqexternalref{::type scope:x '->' x}{http://rocq-prover.org/doc/V9.1.1/stdlib//Corelib.Init.Logic}{\coqdocnotation{\ensuremath{\rightarrow}}} \coqref{SECF.Maps.m':76}{\coqdocvariable{m'}} \coqref{SECF.Maps.x:77}{\coqdocvariable{x}} \coqexternalref{::type scope:x '=' x}{http://rocq-prover.org/doc/V9.1.1/stdlib//Corelib.Init.Logic}{\coqdocnotation{=}} \coqexternalref{Some}{http://rocq-prover.org/doc/V9.1.1/stdlib//Corelib.Init.Datatypes}{\coqdocconstructor{Some}} \coqref{SECF.Maps.v:78}{\coqdocvariable{v}}.\coqdoceol
\coqdocemptyline
\end{coqdoccode}
We can then show that map update preserves map inclusion, that is: 
\begin{coqdoccode}
\coqdocemptyline
\coqdocnoindent
\coqdockw{Lemma} \coqdef{SECF.Maps.includedin update}{includedin\_update}{\coqdoclemma{includedin\_update}} : \coqdockw{\ensuremath{\forall}} (\coqdef{SECF.Maps.A:79}{A}{\coqdocbinder{A}} : \coqdockw{Type}) (\coqdef{SECF.Maps.m:80}{m}{\coqdocbinder{m}} \coqdef{SECF.Maps.m':81}{m'}{\coqdocbinder{m'}} : \coqref{SECF.Maps.partial map}{\coqdocdefinition{partial\_map}} \coqref{SECF.Maps.A:79}{\coqdocvariable{A}})\coqdoceol
\coqdocindent{16.50em}
(\coqdef{SECF.Maps.x:82}{x}{\coqdocbinder{x}} : \coqdocinductive{string}) (\coqdef{SECF.Maps.vx:83}{vx}{\coqdocbinder{vx}} : \coqref{SECF.Maps.A:79}{\coqdocvariable{A}}),\coqdoceol
\coqdocindent{1.00em}
\coqref{SECF.Maps.includedin}{\coqdocdefinition{includedin}} \coqref{SECF.Maps.m:80}{\coqdocvariable{m}} \coqref{SECF.Maps.m':81}{\coqdocvariable{m'}} \coqexternalref{::type scope:x '->' x}{http://rocq-prover.org/doc/V9.1.1/stdlib//Corelib.Init.Logic}{\coqdocnotation{\ensuremath{\rightarrow}}}\coqdoceol
\coqdocindent{1.00em}
\coqref{SECF.Maps.includedin}{\coqdocdefinition{includedin}} (\coqref{SECF.Maps.x:82}{\coqdocvariable{x}} \coqref{SECF.Maps.:::x '|->' x ';' x}{\coqdocnotation{|->}} \coqref{SECF.Maps.vx:83}{\coqdocvariable{vx}} \coqref{SECF.Maps.:::x '|->' x ';' x}{\coqdocnotation{;}} \coqref{SECF.Maps.m:80}{\coqdocvariable{m}}) (\coqref{SECF.Maps.x:82}{\coqdocvariable{x}} \coqref{SECF.Maps.:::x '|->' x ';' x}{\coqdocnotation{|->}} \coqref{SECF.Maps.vx:83}{\coqdocvariable{vx}} \coqref{SECF.Maps.:::x '|->' x ';' x}{\coqdocnotation{;}} \coqref{SECF.Maps.m':81}{\coqdocvariable{m'}}).\coqdoceol
\coqdocnoindent
\coqdockw{Proof}.\coqdoceol
\coqdocindent{1.00em}
\coqdoctac{unfold} \coqref{SECF.Maps.includedin}{\coqdocdefinition{includedin}}.\coqdoceol
\coqdocindent{1.00em}
\coqdoctac{intros} \coqdocvar{A} \coqdocvar{m} \coqdocvar{m'} \coqdocvar{x} \coqdocvar{vx} \coqdocvar{H}.\coqdoceol
\coqdocindent{1.00em}
\coqdoctac{intros} \coqdocvar{y} \coqdocvar{vy}.\coqdoceol
\coqdocindent{1.00em}
\coqdoctac{destruct} (\coqdoclemma{eqb\_spec} \coqdocvar{x} \coqdocvar{y}) \coqdockw{as} [\coqdocvar{Hxy} \ensuremath{|} \coqdocvar{Hxy}].\coqdoceol
\coqdocindent{1.00em}
- \coqdoctac{rewrite} \coqdocvar{Hxy}.\coqdoceol
\coqdocindent{2.00em}
\coqdoctac{rewrite} \coqref{SECF.Maps.update eq}{\coqdoclemma{update\_eq}}. \coqdoctac{rewrite} \coqref{SECF.Maps.update eq}{\coqdoclemma{update\_eq}}. \coqdoctac{intro} \coqdocvar{H1}. \coqdoctac{apply} \coqdocvar{H1}.\coqdoceol
\coqdocindent{1.00em}
- \coqdoctac{rewrite} \coqref{SECF.Maps.update neq}{\coqdoclemma{update\_neq}}.\coqdoceol
\coqdocindent{2.00em}
+ \coqdoctac{rewrite} \coqref{SECF.Maps.update neq}{\coqdoclemma{update\_neq}}.\coqdoceol
\coqdocindent{3.00em}
\ensuremath{\times} \coqdoctac{apply} \coqdocvar{H}.\coqdoceol
\coqdocindent{3.00em}
\ensuremath{\times} \coqdoctac{apply} \coqdocvar{Hxy}.\coqdoceol
\coqdocindent{2.00em}
+ \coqdoctac{apply} \coqdocvar{Hxy}.\coqdoceol
\coqdocnoindent
\coqdockw{Qed}.\coqdoceol
\coqdocemptyline
\end{coqdoccode}
This property is quite useful for reasoning about languages with
    variable binding -- e.g., the Simply Typed Lambda Calculus, which
    we will see in \textit{Programming Language Foundations}, where maps are
    used to keep track of which program variables are defined in a
    given scope. 
\begin{coqdoccode}
\coqdocemptyline
\end{coqdoccode}
\end{document}
